\chapter*{Abstract}
\addcontentsline{toc}{chapter}{Abstract}

This dissertation studies the machine-checked verification of the formal
provable-security model of the lattice-based signature scheme \haetae{}.  The
verification is performed in \easycrypt{} and targets the formal,
reduction-style random-oracle security argument rather than implementation
certification artifacts or a standard-model proof for SHAKE.  The central
machine-checked statement is a scoped \eufcma{} implication for the modeled
scheme in the random oracle model, with the counted-ROM loss ledger
instantiated at the checked wrapper constants and under the explicit hop and
hardness premises of the final EasyCrypt composition lemma.  The current
dissertation also records a vacuity warning: one final
composition premise contains \ec{rejection\_sampling\_loss\_term}, which is
larger than one under the present constants, so the checked implication should
not be read as a non-vacuous concrete security instantiation until that
antecedent or loss normalization is repaired.  It is not claimed as a checked
equivalence theorem between the February 2026 HAETAE specification and the
EasyCrypt model.  The proof surface is supported by formal
modules for signature games, lazy random oracles, transcript games, signing
simulators, paper-simulation samplers, exact hyperball sampling, rejection
sampling, special soundness, and concrete loss accounting.

A major practical contribution is the creation of a separated proof-management
surface under \code{provable-security/}.  The manifest
\code{provable-security/proof-files.txt} lists the EasyCrypt files that define
the provable-security gate, and \code{provable-security/easycrypt/} contains the
corresponding source copies.  The verification script compiles this separated
folder with EasyCrypt and excludes generated known-answer-test material such as
\code{HAETAE\_FIPS202\_KAT\_Empty\_Generated.ec}, which is useful for
implementation validation but not necessary for the reduction proof.

The dissertation explains the theorem architecture from the generic signature
\rom{} interface in \code{Sig\_ROM.ec} through the \haetae{}-specific theorem
layer in \code{HAETAE\_Security.ec}.  It gives particular attention to the
proof-engineering obligations that are easy to hide in informal cryptographic
papers: random-oracle programming freshness, adversarial query logging,
\code{sampler\_bad\_prequery}, clean-conditioned lifting through signing
oracles, exact hyperball sampling, rejection-sampling losses, and the algebra of
final concrete bounds.  It also records a dependency layer separating checked
EasyCrypt theorems, explicit modeling assumptions, audited specification
correspondences, vacuity checks, and out-of-scope cryptanalytic or
implementation claims.  The resulting document is both a dissertation-style
explanation of the checked formal-model proof surface and a guide to maintaining
it.  The model/specification correspondence, non-vacuous discharge of the final
hop premise, concrete hash instantiation, implementation correctness, concrete
MLWE/Module-SIS security estimates, and reduction runtime overhead are
identified as explicit boundary obligations rather than hidden theorem claims.

\cleardoublepage
